\section{Specifying the Upgrade: Rescues, Damages and Net}
\label{sec:value}
\label{sec:r1}
\label{sec:r2}

The holdout is pre-registered; the fragility, outcome-class and label-transfer analyses are exploratory and post hoc.

\paragraph{Many rescues, many damages, small net.} On the \tautwo{} development panel (85 tasks $\times$ 20 complete blocks), replacing a W4 agent's first tool-structure-divergent action with B's, while Q continues, changes the task outcome in $M=15.5$--$20.9$\% of units (Figure~\ref{fig:case}b). Rescues and damages are of similar size, $r=6.8$--$11.4$ and $d=7.5$--$10.1$\pp{}, so the net gain $\Gall$ lies between $-1.9$ and $+2.9$\pp{}. Two fixed-pool intervals exclude zero (negative for q4 GPTQ-W4, positive for q8 GPTQ-W4); no new-task interval does. KV8 configurations turn over fewer outcomes ($M=6.5$--$8.1$\%, a descriptive comparison). $\Gall$ is a local replacement effect, not an end-to-end comparison (Appendix~\ref{app:e2e}). % src: analysis/fork/out/dev_all_20261001T070326Z.md
Pooled over the four W4 configurations, $\Gall=+0.7$\pp{} (fixed pool [$-0.3$, $+1.6$]; new tasks [$-2.1$, $+3.5$]) with $M=18.6$\%. Table~\ref{tab:refs} sets this turnover beside two reference interventions, each against B's action on its own population: a second W4 checkpoint (\S\ref{sec:sameprec}) and a key-order reversal that leaves the operation identical (\S\ref{sec:scope}). % src: analysis/fork/out/dev_all_20261001T070326Z.md (W4:all pooled row: G block bootstrap, M); analysis/fork/out/task_fragility_20261005T134836Z.md (Full pool, pool G crossed)

\begin{table}[t]
\centering\small\setlength{\tabcolsep}{3pt}
% src: analysis/fork/out/dev_all_20261001T070326Z.md (cells M); analysis/fork/out/sameprec_20261004T070601Z.md (TL;DR; Full pool pool rows); analysis/fork/out/keyorder_20261005T204349Z.md (M row)
\begin{tabular}{@{}p{3.55cm}lr@{}}
\toprule
Population (denominator) & Intervention & Turnover \\
\midrule
four W4 configurations, 85 tasks $\times$ 20 blocks (units) & B's action & 15.5--20.9 \\
\addlinespace
q4, q8 GPTQ-W4, 85 tasks $\times$ 6 blocks (units) & B's action & 16.8 \\
 & Q$'$'s action & 17.7 \\
\addlinespace
300 key-order-eligible q8 GPTQ-W4 states (states) & B's action & 9.0 \\
 & key-order reversal & 6.7 \\
\bottomrule
\end{tabular}
\caption{Reference interventions (\%; \S\ref{sec:sameprec}, \S\ref{sec:scope}). Each block of rows is a different population with its own denominator. The references are compared with B's action only within a population; they are not subtracted to obtain a precision-specific turnover.}
\label{tab:refs}
\end{table}

\paragraph{Pre-registered telecom holdout.} On 34 unseen telecom tasks, under a protocol frozen before any holdout unit ran (12 fresh blocks, three W4 configurations), the replacement again rescued and damaged ($r=5.2$, $d=9.8$\pp{}). The primary $\Gall=-4.57$\pp{} (crossed 90\% [$-7.17$, $-1.96$]) lies below zero, but the frozen readings had only positive-gain and small-effect categories, so the protocol classifies it as \emph{inconclusive} (Appendix~\ref{app:prereg}). Both pre-registered secondary checks passed: S1 (rescues and damages both positive) and S2 (cross-half repeatability of task fragility, $\rho=0.588$, one-sided permutation $p=0.0001$). All 135 QQ no-op replays selected by the fixed 10\% hash rule matched. % src: analysis/fork/out/holdout_confirm_20261002T012353Z.md (TL;DR; Readings; S1; S2); docs/insights/PREREG_FORK_CONFIRM.md §0, §8; analysis/fork/out/holdout_noop_audit_20261004T152207Z.md (TL;DR)
In a post hoc four-branch analysis of the holdout's divergent states, the pooled $H$ and $F$ are also negative ($-2.8$ and $-2.5$\pp{}), but both crossed 95\% intervals include zero, and q4 GPTQ-W4 has positive point estimates (Appendix Table~\ref{tab:holdoutr2}). The holdout changes domain as well as tasks; on the 80 telecom development tasks the same q8 configurations give $+2.7$ and $-0.6$\pp{} (Appendix~\ref{app:holdoutdev}). % src: analysis/fork/out/holdout_r2_posthoc_20261002T024414Z.md (§1 crossed 95%); analysis/fork/out/dev_all_20261001T093944Z.md (Telecom development)

\paragraph{The consequences recur on the same tasks.} We split the 20 nonce blocks into two halves by an outcome-independent hash of the block name, ranked tasks by turnover on one half and measured the other half (Appendix Table~\ref{tab:fragility}). The 17 most fragile tasks (top 20\%) carry 41.0\% (crossed [26.6, 55.9]) and 42.3\% ([27.0, 57.5]) of the held-out turnover in the two directions, against 20\% if fragility did not persist. Their turnover is 38.7\%, while the other 68 tasks still turn over 13.7\%. Both groups have rescues and damages; neither group's net gain is resolved under the crossed interval. This replicates fragility on other nonces of the same tasks; it does not predict which new tasks are fragile. % src: analysis/fork/out/task_fragility_20261005T134836Z.md (Split-half fragility, both directions); analysis/fork/out/task_fragility_split_20261005T134815Z.json; derived: 68=85-17

\paragraph{Which mode is evaluated decides which rescues are counted.} On divergent states, the average effect of the action swap has the same point estimate under Q's and B's continuation (W4 pool: $\DQ=\DB=+0.5$\pp{}), and no average gain is detected for upgrading this step, the rest, or both ($G=+0.5$, $H=+0.2$, $F=+0.7$\pp{}; all fixed-pool and new-task intervals cover zero, Appendix~\ref{app:r2main}). % src: analysis/fork/out/ghf_20261001T211519Z.md (§1 pools, W4:all; F - H = Δ_B); analysis/fork/out/ghf_20261001T211519Z.json
Averages hide states where the single-mode evaluations disagree. In a \emph{joint rescue}, only $\YY{BB}$ succeeds (Figure~\ref{fig:case}a). These are 4.0\% of W4 divergent states (fixed-pool 95\% [3.5, 4.5]; new-task [2.7, 5.3]); testing only the current action misses them. In a \emph{rescue undone}, B's action rescues under Q's continuation ($\YY{BQ}=1$, $\YY{QQ}=0$) but B's takeover undoes it ($\YY{BB}=0$). These are 4.8\% (fixed-pool [4.1, 5.4]; new-task [3.2, 6.5]); testing only a whole takeover misses them. % src: analysis/fork/out/outcome_classes_20261004T143614Z.md
A conditional label-permutation reference (which keeps how many of the four branches succeed) already produces both patterns at similar size (joint rescue 4.0\% vs 3.5\%; rescue undone 4.8\% vs 5.8\%), so the frequencies give no generalisable evidence of extra enrichment (Appendix~\ref{app:offdiag}); the reference does not model execution, so it does not say how these outcomes arise. Whatever their cause, the patterns mean that a gain measured for action replacement cannot be reused as the gain of a whole takeover, or the reverse. % src: analysis/fork/out/pattern_ref_20261005T134721Z.md (§1, pool W4:all)

\paragraph{Intervention-label transfer.} We rank tasks by one mode's mean gain on half of the nonce blocks and execute a mode on the same tasks in the other half (known tasks across nonces, not new tasks or single decisions). The frozen post hoc protocol designates ordinary states (Bank U, U1) as the primary frame and first divergences (Bank P) as secondary; it fixes no budget or label pair, and we report one selected entry, the continuation label choosing joint replacement at a 10\% task budget (full matrix in Appendix~\ref{app:transfer}; crossed 95\%, uncorrected). In the primary frame, selecting tasks for joint replacement by the continuation label instead of the joint label loses $+0.30$\pp{} [$-0.60$, $+1.11$], not resolved. In the secondary frame, the joint-replacement label gains $+4.35$\pp{} of all units and the continuation label $+1.52$\pp{}, a paired loss of $2.82$\pp{} [1.40, 3.90]. % src: analysis/fork/out/label_transfer_20261005T185537Z.md (Bank U U1 PRIMARY split-half budget 10%, row H, column F: loss; Bank P dev_all units SECONDARY split-half budget 10%); docs/insights/LABEL_TRANSFER_POSTHOC.md (§2 frames)
